\documentclass[14pt, a4paper]{extarticle}
\usepackage[T2A]{fontenc}
\usepackage[utf8]{inputenc}
\usepackage[english, russian]{babel}
\usepackage[left=30mm, right=10mm, top=20mm, bottom=25mm]{geometry}
\usepackage{amsmath, amssymb, amsfonts}
\usepackage{graphicx}
\usepackage{float}
\usepackage{hyperref}
\usepackage{longtable}
\usepackage{booktabs}
\usepackage{multirow}
\usepackage{enumitem}
\usepackage{array}
\usepackage{setspace}

\usepackage{misccorr}
\usepackage{indentfirst}

\usepackage{tikz}
\usetikzlibrary{arrows.meta, positioning, shapes.geometric}
\usepackage{listings}
\usepackage{xcolor}
\usepackage{caption}

\setstretch{1.2}

% Заголовки section по центру
\makeatletter
\renewcommand\section{\@startsection {section}{1}{\z@}%
{-3.5ex \@plus -1ex \@minus -.2ex}%
{2.3ex \@plus.2ex}%
{\centering\normalfont\Large\bfseries}}
\makeatother

% Точки в списке литературы
\makeatletter
\renewcommand{\@biblabel}{#1.}
\makeatother

% Настройка hyperref
\hypersetup{
colorlinks=true,
linkcolor=black,
citecolor=black,
urlcolor=blue
}

\lstdefinestyle{portalcode}{
basicstyle=\ttfamily\small,
columns=fullflexible,
breaklines=true,
frame=single,
numbers=left,
numberstyle=\tiny,
keywordstyle=\bfseries,
commentstyle=\itshape,
showstringspaces=false
}

\begin{document}

% Титульная страница
\thispagestyle{empty}
{\scriptsize
\begin{center}
ФЕДЕРАЛЬНОЕ ГОСУДАРСТВЕННОЕ ОБРАЗОВАТЕЛЬНОЕ БЮДЖЕТНОЕ УЧРЕЖДЕНИЕ\\
ВЫСШЕГО ОБРАЗОВАНИЯ
\end{center}
}
{\bf
\begin{center}
ФИНАНСОВЫЙ УНИВЕРСИТЕТ\\
ПРИ ПРАВИТЕЛЬСТВЕ РОССИЙСКОЙ ФЕДЕРАЦИИ\\
Липецкий филиал
\end{center}
}
\hrule
\bigskip

{\small
\begin{center}
\textbf{
Кафедра <<Учет и информационные технологии в бизнесе>>}
\end{center}
}
\bigskip

Направление: Прикладная математика и информатика

\medskip
Профиль: Анализ данных и принятие решений в экономике и финансах

\vskip 2cm

\begin{center}
\textbf{КУРСОВОЙ ПРОЕКТ}
\end{center}

по дисциплине: <<Машинное обучение>>

\bigskip
ТЕМА: Проектирование и разработка прототипа интерактивного
образовательного веб-портала для пациентов с впервые выявленным
сахарным диабетом 1 типа

\vskip 3cm

\hskip .3\textwidth Преподаватель: Калитвин Владимир Анатольевич

\hskip .3\textwidth Студент: Кузнецов Никита Владимирович

\hskip .3\textwidth Учетный номер \textnumero 100.16/240075

\hskip .3\textwidth Курс второй

\vskip 4cm

\begin{center}
Липецк 2026
\end{center}

\newpage

% Содержание
\tableofcontents
\newpage

\addcontentsline{toc}{section}{ВВЕДЕНИЕ}
\section*{ВВЕДЕНИЕ}

\noindent Сахарный диабет 1 типа (СД1) — это хроническое аутоиммунное заболевание, при котором поджелудочная железа перестаёт вырабатывать инсулин. В России, по данным Федерального регистра сахарного диабета на 2025 год, зарегистрировано более 350 тысяч пациентов с СД1, и ежегодно выявляется порядка 15–20 тысяч новых случаев. Значительную долю среди заболевших составляют дети, подростки и молодые люди до 30 лет. Диагноз «сахарный диабет 1 типа» звучит как приговор, но на самом деле это не конец жизни, а начало новой — с постоянным контролем, подсчётами и самодисциплиной.

\medskip
\noindent Когда мне предложили выбрать тему для курсового проекта, я сразу подумал о чём-то социально значимом. Хотелось не просто написать код «для галочки», а сделать что-то, что реально может помочь людям. Тема диабета мне близка: у моего знакомого в 19 лет выявили СД1, и я видел, как тяжело ему было в первые месяцы — куча вопросов, паника, непонимание, что делать. Врачи в поликлинике дают общие рекомендации, а дальше человек остаётся один на один с болезнью. В интернете информации много, но она разрозненная, часто противоречивая, а порой и откровенно вредная. Именно так родилась идея создать образовательный веб-портал, который поможет пациентам с впервые выявленным СД1 разобраться в своём заболевании.

\medskip
\noindent Современные веб-технологии позволяют сделать обучение не скучным чтением простыней текста, а интерактивным процессом с калькуляторами, дневниками, тестами и визуализацией данных. Машинное обучение здесь тоже может найти применение: например, для прогнозирования уровня глюкозы на основе предыдущих измерений или для персонализированных рекомендаций по питанию. Но в рамках данного курсового проекта мы сосредоточимся на проектировании и разработке прототипа, а внедрение ML-алгоритмов оставим на будущее.
\vskip 5cm
\subsection*{Актуальность темы}

\noindent Актуальность разработки образовательного веб-портала для пациентов с СД1 обусловлена несколькими факторами. Во-первых, число заболевших неуклонно растёт. По прогнозам Международной диабетической федерации (IDF), к 2030 году количество пациентов с СД1 в мире увеличится на 20–25\%. 

\medskip 
\noindent Во-вторых, качество жизни пациента напрямую зависит от его осведомлённости о болезни. Исследования показывают, что пациенты, прошедшие структурированное обучение, имеют уровень гликированного гемоглобина (HbA1c) в среднем на 1–1,5\% ниже, чем те, кто не обучался. Каждый процент снижения HbA1c — это минус 30–35\% риска осложнений: ретинопатии, нефропатии, нейропатии.

\medskip
\noindent В-третьих, существующие образовательные ресурсы либо слишком сложны для новичков (написаны медицинским языком с кучей терминов), либо слишком примитивны (дают общие советы без учёта индивидуальных особенностей). Нет единой платформы, которая бы сочетала в себе теорию, практические инструменты и возможность отслеживать свой прогресс. Именно такую платформу мы и попытаемся создать.

\subsection*{Цель и задачи работы}

\noindent \textbf{Целью} данной курсовой работы является проектирование и разработка прототипа интерактивного образовательного веб-портала для пациентов с впервые выявленным сахарным диабетом 1 типа.

\medskip
\noindent Для достижения поставленной цели необходимо решить следующие \textbf{задачи}:

\begin{enumerate}
\item Изучить теоретические основы сахарного диабета 1 типа: этиологию, патогенез, принципы лечения и роль терапевтического обучения пациентов.

\item Провести обзор существующих образовательных ресурсов для пациентов с СД1, выявить их достоинства и недостатки.

\item Спроектировать структуру веб-портала, определить функциональные модули и сценарии использования.

\item Разработать математическую модель расчёта хлебных единиц как ключевого элемента планирования питания при СД1.

\item Реализовать прототип клиентской части веб-портала с использованием современных веб-технологий.

\item Описать алгоритмы работы интерактивных компонентов: калькулятора ХЕ, дневника самочувствия, системы тестирования.

\item Разработать структуру базы данных для хранения пользовательской информации и результатов измерений.

\item Оформить техническую документацию в соответствии с требованиями ГОСТ.
\end{enumerate}

\subsection*{Объект и предмет исследования}

\noindent \textbf{Объектом исследования} является процесс обучения пациентов с впервые выявленным сахарным диабетом 1 типа.

\medskip
\noindent \textbf{Предметом исследования} выступает интерактивный образовательный веб-портал как инструмент поддержки пациентов на начальном этапе адаптации к заболеванию.

\subsection*{Методы исследования}

\noindent При выполнении работы использовались следующие методы: анализ научно-медицинской литературы и клинических рекомендаций, сравнительный анализ существующих образовательных платформ, метод прототипирования при проектировании интерфейса, методы математического моделирования при разработке калькулятора хлебных единиц, экспериментальное тестирование работоспособности прототипа.
\vskip 5cm
\subsection*{Структура работы}

\noindent Курсовая работа состоит из введения, трёх глав, заключения, списка литературы и приложений. Общий объём работы составляет 33 страниц машинописного текста.

\medskip
\noindent В первой главе рассматриваются теоретические основы сахарного диабета 1 типа и современные подходы к терапевтическому обучению пациентов. Описывается роль структурированных образовательных программ в достижении компенсации углеводного обмена.

\medskip
\noindent Во второй главе проводится анализ существующих решений для пациентов с СД1, обосновывается выбор технологического стека для разработки прототипа, формулируются функциональные требования к системе.

\medskip
\noindent Третья глава посвящена проектированию и реализации прототипа веб-портала: описывается структурная схема, математическая модель расчёта хлебных единиц, алгоритмы работы модулей, структура базы данных, макеты интерфейса.

\newpage
\section{ТЕОРЕТИЧЕСКИЕ ОСНОВЫ САХАРНОГО ДИАБЕТА 1 ТИПА И ОБУЧЕНИЯ ПАЦИЕНТОВ}

\subsection{Сахарный диабет 1 типа: этиология, патогенез и клиническая картина}

\noindent Сахарный диабет 1 типа — это хроническое заболевание, в основе которого лежит аутоиммунное разрушение бета-клеток поджелудочной железы, вырабатывающих инсулин. Без инсулина глюкоза не может проникнуть в клетки, и организм буквально «голодает» при высоком уровне сахара в крови. Если совсем просто: представьте, что клетка — это дом, глюкоза — это дрова для печки, а инсулин — ключ от двери. При СД1 ключа нет, дрова лежат на улице, а в доме холодно.

\medskip
\noindent Этиология СД1 до конца не изучена, но известно, что заболевание возникает при сочетании генетической предрасположенности и провоцирующих факторов внешней среды. Гены тут играют важную, но не определяющую роль. Если у одного из родителей СД1, риск для ребёнка составляет 5–8\%, если оба родителя больны — до 20\%. Среди провоцирующих факторов выделяют вирусные инфекции (энтеровирусы, краснуха, эпидемический паротит), особенности питания в раннем детстве (раннее введение коровьего молока), дефицит витамина D.

\medskip
\noindent Патогенез СД1 можно разделить на несколько стадий. На первой стадии у генетически предрасположенного человека появляются аутоантитела к бета-клеткам (к островковым клеткам, к инсулину, к глутаматдекарбоксилазе — GAD). На этом этапе никаких симптомов нет, уровень глюкозы в норме. На второй стадии начинается постепенное разрушение бета-клеток, появляются первые нарушения углеводного обмена — повышение глюкозы после еды. На третьей стадии, когда погибает около 80–90\% бета-клеток, развивается клиническая картина: жажда, частое мочеиспускание, потеря веса, слабость, запах ацетона изо рта. Если вовремя не начать лечение, развивается диабетический кетоацидоз — угрожающее жизни состояние, требующее немедленной госпитализации.

\vskip  1cm
\noindent Клиническая картина СД1 развивается, как правило, остро — в течение нескольких недель, реже месяцев. Пациент начинает много пить (полидипсия), много мочиться (полиурия), быстро теряет вес при нормальном или даже повышенном аппетите. Характерный симптом — слабость, сонливость, снижение работоспособности. У детей может наблюдаться ночное недержание мочи, которого раньше не было. Если появились такие симптомы — нужно срочно сдать кровь на сахар, и если он выше 11 ммоль/л в любое время суток — это практически гарантированный диабет.

\medskip
\noindent Лечение СД1 заключается в пожизненной заместительной инсулинотерапии. Инсулин вводится подкожно с помощью шприц-ручек или инсулиновой помпы. Доза подбирается индивидуально и зависит от многих факторов: веса, возраста, физической активности, питания, сопутствующих заболеваний. В среднем суточная потребность в инсулине составляет 0,5–0,8 ЕД на килограмм веса, но это очень приблизительная цифра.

\subsection{Роль терапевтического обучения в управлении сахарным диабетом}

\noindent Терапевтическое обучение пациентов (ТОП) — это не просто информирование о болезни, а структурированный процесс, направленный на формирование у пациента навыков самостоятельного управления заболеванием. Всемирная организация здравоохранения признала ТОП неотъемлемой частью лечения хронических заболеваний ещё в 1998 году, и с тех пор эта концепция только укрепилась.

\medskip
\noindent Почему обучение так важно именно при СД1? Потому что ни один врач не может круглосуточно находиться рядом с пациентом и говорить ему: «сейчас тебе нужно съесть столько-то, уколоть столько-то, пойти погулять столько-то минут». Все эти решения пациент принимает самостоятельно каждый день, по 10–15 раз. И от правильности этих решений зависит, разовьются ли у него через 10–15 лет осложнения или нет.

\vskip 5cm
\noindent Исследования показывают впечатляющие результаты. Знаменитое исследование DCCT (Diabetes Control and Complications Trial), которое проводилось в США и Канаде с 1983 по 1993 год, доказало: интенсивное лечение и обучение пациентов снижает риск развития диабетической ретинопатии на 76\%, нейропатии на 60\%, нефропатии на 50\%. Эти цифры говорят сами за себя. Последующее наблюдение за участниками исследования (EDIC) показало, что эффект сохраняется десятилетиями — феномен, который назвали «метаболической памятью».

\medskip
\noindent В России терапевтическое обучение пациентов с СД1 регламентируется клиническими рекомендациями, утверждёнными Минздравом. Обучение должно проводиться по структурированным программам, которые включают несколько обязательных модулей: что такое диабет и чем он опасен, самоконтроль глюкозы крови, принципы инсулинотерапии, питание и подсчёт хлебных единиц, физическая активность и её коррекция, профилактика и лечение гипогликемии, диабет и сопутствующие заболевания, психологические аспекты жизни с диабетом.

\medskip
\noindent Однако на практике ситуация далека от идеальной. В поликлиниках эндокринологи физически не успевают проводить полноценное обучение — на приём одного пациента отводится 12–15 минут. За это время можно только выписать рецепт и спросить, как дела. Школы диабета, которые раньше работали при многих больницах, сокращаются из-за нехватки финансирования. Пациенты из малых городов и сёл вообще лишены возможности пройти структурированное обучение очно. Именно здесь на помощь могут прийти цифровые технологии.
\vskip 7cm
\subsection{Современные подходы к дистанционному обучению пациентов}

\noindent Развитие интернет-технологий и мобильных приложений открыло новые возможности для обучения пациентов. Дистанционное обучение имеет ряд преимуществ перед очным: доступность 24/7 из любой точки, где есть интернет; возможность учиться в своём темпе, возвращаться к сложным темам; интерактивность — тесты, калькуляторы, визуализация данных; экономическая эффективность — один раз созданный курс могут проходить тысячи пациентов.

\noindent Существует несколько форматов дистанционного обучения: вебинары и онлайн-школы в реальном времени; записанные видеолекции; структурированные текстовые курсы с тестами для самопроверки; мобильные приложения с дневниками самоконтроля и калькуляторами; чат-боты, имитирующие диалог с консультантом.

\medskip
\noindent Наиболее эффективным считается смешанный подход: очное первичное консультирование плюс дистанционное сопровождение. Но если очное обучение недоступно (а так бывает часто), то хорошо сделанный образовательный портал может хотя бы частично его заменить.

\medskip
\noindent В России тема дистанционного обучения пациентов с СД1 развита недостаточно. Есть несколько ресурсов: сайт «Диабет.ру» с базой статей, форум «Диа-клуб», YouTube-канал «Диабет — образ жизни», отдельные мобильные приложения. Но все они решают лишь часть задач: где-то есть теория, но нет калькулятора; где-то есть калькулятор, но нет дневника; где-то есть дневник, но нет системы тестирования. Комплексного решения, объединяющего все необходимые инструменты на одной платформе, на данный момент не существует.

\newpage
\section{АНАЛИЗ ПРЕДМЕТНОЙ ОБЛАСТИ И ПОСТАНОВКА ЗАДАЧИ}

\subsection{Обзор существующих образовательных ресурсов для пациентов с СД1}

\noindent Прежде чем приступить к проектированию собственного решения, был проведён анализ существующих ресурсов. Хотелось понять, что уже есть, какие функции востребованы, а какие — отсутствуют. Анализ проводился по следующим критериям: полнота теоретической базы, наличие интерактивных инструментов (калькуляторы, дневники), удобство интерфейса, актуальность информации, стоимость использования.

\subsubsection{Сайт «Диабет.ру»}

\noindent Один из старейших русскоязычных ресурсов о диабете. Содержит большую базу статей по всем аспектам заболевания: от причин и симптомов до рецептов блюд. Есть форум, где пациенты общаются между собой.

\medskip
\noindent \textbf{Достоинства:} много информации, активное сообщество, бесплатный доступ.

\medskip
\noindent \textbf{Недостатки:} дизайн устаревший (как будто из 2008 года), статьи не структурированы в учебные модули, нет интерактивных инструментов — только текст. Новичку сложно понять, с чего начать, в каком порядке читать статьи.

\subsubsection{Приложение «MySugr»}

\noindent Австрийское приложение для ведения дневника самоконтроля. Позволяет записывать уровень глюкозы, дозы инсулина, количество хлебных единиц, физическую активность. Имеет gamification-элементы: за регулярное ведение дневника пользователь получает баллы и «укрощает диа-монстра».

\medskip
\noindent \textbf{Достоинства:} отличный интерфейс, мотивация к ведению дневника, интеграция с глюкометрами через Bluetooth.

\medskip
\noindent \textbf{Недостатки:} нет образовательного контента, только дневник. Бесплатная версия имеет ограниченный функционал, полная стоит 3–5 евро в месяц. Нет полноценной русской локализации.

\subsubsection{Приложение «Diabetes:M»}

\noindent Ещё одно популярное приложение-дневник. Имеет встроенную базу продуктов с содержанием углеводов, калькулятор ХЕ, возможность построения графиков и экспорта данных в PDF для врача.

\medskip
\noindent \textbf{Достоинства:} мощный функционал, подробная статистика, экспорт отчётов.

\medskip
\noindent \textbf{Недостатки:} интерфейс перегружен, новому пользователю сложно разобраться. Образовательного контента нет. Полная версия — около 20 долларов.

\subsubsection{Школа диабета на YouTube}

\noindent На YouTube есть несколько каналов, посвящённых диабету, где врачи-эндокринологи читают лекции, разбирают реальные случаи, отвечают на вопросы.

\medskip
\noindent \textbf{Достоинства:} бесплатно, информация от практикующих врачей, наглядно.

\medskip
\noindent \textbf{Недостатки:} нет структуры — видео выходят в разном порядке, не связаны в единый курс. Нет возможности проверить усвоение материала. Не всегда удобно смотреть видео (в транспорте, на работе).

\medskip
\noindent \textbf{Вывод по обзору:} существующие решения либо предоставляют только теоретическую информацию, либо только инструменты для самоконтроля. Комплексной платформы, которая сочетала бы образовательные модули, интерактивные инструменты и систему проверки знаний, на данный момент нет. Именно эту нишу должен заполнить разрабатываемый веб-портал.

\subsection{Формирование требований к образовательному веб-порталу}

\noindent На основе анализа предметной области, изучения клинических рекомендаций и существующих решений были сформулированы требования к разрабатываемому прототипу. Требования разделены на функциональные и нефункциональные.
\vskip 5cm
\subsubsection{Функциональные требования}

\begin{enumerate}
\item \textbf{Личный кабинет пользователя.} Система должна позволять пользователю зарегистрироваться, ввести основные данные (имя, email, дату постановки диагноза, тип инсулина), хранить историю прохождения образовательных модулей и результаты тестирования.

\item \textbf{Образовательные модули.} Система должна содержать структурированный теоретический материал по следующим темам: «Что такое диабет?», «Питание и хлебные единицы», «Физическая активность», «Самоконтроль и инсулинотерапия». Каждый модуль должен завершаться проверочным тестом.

\item \textbf{Калькулятор хлебных единиц.} Пользователь должен иметь возможность выбрать продукт из базы, указать массу порции и получить расчёт количества ХЕ. Также должна быть возможность сложить несколько продуктов и получить суммарное количество ХЕ за приём пищи.

\item \textbf{Дневник самочувствия.} Пользователь должен иметь возможность записывать: дату и время, уровень глюкозы до и после еды, количество съеденных ХЕ, дозу инсулина, физическую активность, субъективные ощущения (хорошо, слабость, гипогликемия).

\item \textbf{Система тестирования.} После каждого образовательного модуля пользователю предлагается пройти тест из 8–12 вопросов. Система показывает результат в процентах и даёт рекомендации.

\item \textbf{Визуализация данных.} На основе записей дневника система строит графики изменения уровня глюкозы, отображает статистику (среднее значение, вариабельность, процент времени в целевом диапазоне).

\item \textbf{Административная панель.} Для управления контентом: добавление и редактирование образовательных материалов, управление пользователями, просмотр статистики использования.
\end{enumerate}

\subsubsection{Нефункциональные требования}

\begin{enumerate}
\item \textbf{Доступность.} Портал должен корректно отображаться на устройствах с разным разрешением экрана (компьютер, планшет, смартфон).

\item \textbf{Простота использования.} Интерфейс должен быть интуитивно понятным для пользователей без технической подготовки. Все надписи и сообщения — на русском языке, без медицинского жаргона.

\item \textbf{Безопасность.} Данные пользователей должны храниться в защищённом виде, доступ к личному кабинету — по паролю.

\item \textbf{Производительность.} Время загрузки страниц не должно превышать 2–3 секунды. Расчёт ХЕ должен происходить мгновенно.

\item \textbf{Масштабируемость.} Архитектура системы должна позволять добавление новых модулей и функций без полной переработки.
\end{enumerate}

\subsection{Выбор технологического стека}

\noindent Для реализации прототипа необходимо выбрать технологии, которые позволят создать работающий продукт в рамках ограниченного времени (один семестр) и с учётом моих текущих знаний в программировании.

\medskip
\noindent \textbf{Язык программирования:} Python. Выбор Python обусловлен несколькими причинами: во-первых, это основной язык, который мы изучаем на направлении «Прикладная информатика»; во-вторых, Python имеет огромное количество библиотек для веб-разработки, работы с данными и машинного обучения; в-третьих, на Python можно быстро создать прототип и при необходимости легко его доработать.

\medskip
\noindent \textbf{Фреймворк для веб-интерфейса:} Streamlit. Почему не Flask или Django, которые обычно используют для веб-разработки? Потому что Streamlit позволяет создавать веб-интерфейсы на чистом Python, без необходимости писать HTML, CSS и JavaScript. Для прототипа, где важна скорость разработки и не требуется сложная серверная логика, Streamlit подходит идеально. Он, конечно, не заменит полноценный фреймворк для коммерческого продукта, но для курсового проекта — самое то.

\vskip 5cm
\noindent \textbf{База данных:} SQLite. Для прототипа не нужна мощная СУБД типа PostgreSQL или MySQL. SQLite — это легковесная встраиваемая база данных, которая хранится в одном файле. Её не нужно настраивать, она работает «из коробки».

\medskip
\noindent \textbf{Дополнительные библиотеки:} для визуализации данных — Plotly (интерактивные графики), для математических расчётов — NumPy (работа с массивами данных, статистические функции), для работы с датами — datetime.

\newpage
\section{ПРОЕКТИРОВАНИЕ И РЕАЛИЗАЦИЯ ПРОТОТИПА ВЕБ-ПОРТАЛА}

\subsection{Структурная схема веб-портала}

\noindent На основе функциональных требований была разработана структурная схема прототипа. Портал состоит из четырёх основных блоков, доступных из личного кабинета: образовательные модули, интерактивные инструменты, аналитика и визуализация, административная панель.

\medskip
\noindent Образовательные модули представляют собой последовательный курс из четырёх тем. Каждая тема содержит теоретический материал и завершается проверочным тестом. Модули расположены в логическом порядке: сначала пациент узнаёт, что такое диабет, потом — как питаться, затем — как заниматься спортом, и наконец — как контролировать уровень глюкозы.

\medskip
\noindent Интерактивные инструменты — это калькулятор хлебных единиц, дневник самочувствия и система тестов. Калькулятор помогает планировать питание, дневник — фиксировать показатели, тесты — проверять знания.

\medskip
\noindent Блок аналитики получает данные из дневника и строит графики, которые помогают увидеть динамику. Например, график уровня глюкозы за неделю или за месяц, диаграмма распределения времени в целевом диапазоне.

\begin{figure}[H]
\centering
\begin{tikzpicture}[
node distance=11mm,
block/.style={rectangle, draw, rounded corners, align=center,
minimum width=38mm, minimum height=10mm},
line/.style={-{Latex[length=3mm]}, thick}
]
\node[block] (account) {Личный кабинет\\пользователя};
\node[block, below left=of account, xshift=-14mm] (modules) {Образовательные\\модули};
\node[block, below=of account] (tools) {Интерактивные\\инструменты};
\node[block, below right=of account, xshift=14mm] (analytics) {Аналитика\\и визуализация};

\node[block, below=of modules] (m1) {Что такое\\диабет?};
\node[block, below=of m1] (m2) {Питание};
\node[block, below=of m2] (m3) {Физическая\\активность};
\node[block, below=of m3] (m4) {Самоконтроль};

\node[block, below=of tools] (calc) {Калькулятор ХЕ};
\node[block, below=of calc] (diary) {Дневник\\самочувствия};
\node[block, below=of diary] (tests) {Проверочные\\тесты};

\draw[line] (account) -- (modules);
\draw[line] (account) -- (tools);
\draw[line] (account) -- (analytics);
\draw[line] (modules) -- (m1);
\draw[line] (m1) -- (m2);
\draw[line] (m2) -- (m3);
\draw[line] (m3) -- (m4);
\draw[line] (tools) -- (calc);
\draw[line] (calc) -- (diary);
\draw[line] (diary) -- (tests);
\draw[line] (diary.east) -- (analytics.south);
\end{tikzpicture}
\caption{Структурная схема прототипа веб-портала}
\end{figure}

\subsection{Математическая модель расчёта хлебных единиц}

\noindent Хлебная единица (ХЕ) — это условная мера, используемая для оценки количества углеводов в продуктах питания. Одна ХЕ соответствует 10–12 граммам углеводов (в разных странах приняты разные стандарты, в России — 12 граммов). Система хлебных единиц была разработана немецкими диетологами специально для пациентов с сахарным диабетом и с тех пор широко используется во всём мире.

\medskip
\noindent Почему это важно? Потому что доза инсулина, которую нужно ввести перед едой, напрямую зависит от количества углеводов, которые пациент планирует съесть. Если пациент знает, сколько ХЕ в его тарелке, он может точно рассчитать дозу инсулина, используя индивидуальный углеводный коэффициент (УК). УК — это количество граммов углеводов, на которое рассчитана 1 единица инсулина. У разных людей УК разный, обычно от 5 до 20 граммов на единицу.

\medskip
\noindent Для каждого продукта количество хлебных единиц определяется по формуле:
\begin{equation}
XEi = \frac{mi \cdot C_i}{100 \cdot 12},
\end{equation}
где $XEi$ — количество хлебных единиц в продукте; $mi$ — масса продукта в граммах; $C_i$ — количество углеводов на 100 граммов продукта; 12 — количество граммов углеводов, соответствующее одной хлебной единице.

\medskip
\noindent Общее количество хлебных единиц в приёме пищи рассчитывается как сумма ХЕ всех продуктов:
\begin{equation}
XE{total}=\sum{i=1}^{n} XE_i.
\end{equation}

\medskip
\noindent Давайте разберём на конкретном примере, чтобы было понятно. Допустим, на завтрак планируется: хлеб пшеничный (60 граммов, углеводов 50 г на 100 г), яблоко (150 граммов, углеводов 12 г на 100 г) и стакан молока (200 граммов, углеводов 5 г на 100 г). Тогда:

\begin{itemize}
\item Для хлеба: XE = (60 × 50) / (100 × 12) = 3000 / 1200 = 2,50 ХЕ.
\item Для яблока: XE = (150 × 12) / (100 × 12) = 1800 / 1200 = 1,50 ХЕ.
\item Для молока: XE = (200 × 5) / (100 × 12) = 1000 / 1200 = 0,83 ХЕ.
\item Итого: 2,50 + 1,50 + 0,83 = 4,83 ХЕ.
\end{itemize}

\noindent Теперь пациент, зная свой углеводный коэффициент, может рассчитать дозу инсулина. Например, если УК = 10 (1 единица инсулина на 10 граммов углеводов), то на 4,83 ХЕ = 4,83 × 12 = 58 граммов углеводов потребуется 58 / 10 = 5,8 единиц инсулина. Округляем до 6 единиц (инсулиновые шприц-ручки обычно имеют шаг 0,5 или 1 единица).

\begin{longtable}{|c|p{0.30\textwidth}|c|c|c|}
\caption{Пример расчёта хлебных единиц для завтрака}\\
\hline
\textbf{№} & \textbf{Продукт} & \textbf{Масса, г} & \textbf{Углеводы, г/100 г} & \textbf{ХЕ} \\
\hline
1 & Хлеб пшеничный & 60 & 50 & 2{,}50 \\
\hline
2 & Яблоко & 150 & 12 & 1{,}50 \\
\hline
3 & Молоко & 200 & 5 & 0{,}83 \\
\hline
\multicolumn{4}{|r|}{\textbf{Итого}} & \textbf{4{,}83} \\
\hline
\end{longtable}

\subsection{Функциональные компоненты прототипа}

\noindent В таблице ниже представлены все функциональные компоненты, которые были запроектированы и частично реализованы в прототипе.

\begin{longtable}{|p{0.24\textwidth}|p{0.34\textwidth}|p{0.30\textwidth}|}
\caption{Функциональные компоненты прототипа веб-портала}\\
\hline
\textbf{Компонент} & \textbf{Назначение} & \textbf{Результат для пользователя} \\
\hline
Личный кабинет & Сохранение данных пользователя, прогресса обучения и результатов тестирования & Индивидуальная траектория обучения и доступ к истории действий \\
\hline
Образовательный модуль & Последовательное изучение теоретического материала по ключевым темам СД1 & Формирование базовых знаний о заболевании \\
\hline
Калькулятор ХЕ & Расчёт количества хлебных единиц на основе выбранных продуктов и массы порции & Упрощение планирования питания \\
\hline
Дневник самочувствия & Фиксация уровня глюкозы, дозы инсулина, активности и субъективных ощущений & Накопление данных для анализа состояния \\
\hline
Система тестов & Проверка усвоения материала после прохождения модулей & Обратная связь и контроль знаний \\
\hline
Визуализация & Отображение графиков, диаграмм и таблиц динамики показателей & Наглядный анализ изменений уровня глюкозы \\
\hline
\end{longtable}
\vskip 5cm
\subsection{Алгоритм проверки данных электронного дневника}

\noindent Одна из частых проблем при ведении дневника — ошибки ввода. Человек может случайно ввести уровень глюкозы 250 ммоль/л вместо 2,5 (перепутать точку и запятую), или указать 500 граммов углеводов за один приём пищи (что физически невозможно). Чтобы такие ошибки не попадали в базу данных и не искажали статистику, необходима проверка вводимых значений.

\medskip
\noindent Алгоритм проверки работает следующим образом. Когда пользователь заполняет форму дневника и нажимает «Сохранить», система сначала проверяет формат введённых данных: все ли обязательные поля заполнены, числа ли в числовых полях, правильный ли формат даты. Если формат неверный — пользователь получает сообщение об ошибке и может исправить данные.

\medskip
\noindent Если формат верный, система проверяет, попадают ли значения в допустимые диапазоны: уровень глюкозы — от 1,0 до 35,0 ммоль/л; количество ХЕ — от 0 до 30 за один приём пищи; доза инсулина — от 0 до 100 единиц; продолжительность активности — от 0 до 300 минут. Если значение выходит за границы диапазона, система предупреждает пользователя, но не блокирует сохранение (вдруг у человека действительно аномальное значение?).

\medskip
\noindent Только после успешного прохождения обеих проверок запись сохраняется в базу данных, и система обновляет графики.


\subsection{Формулы статистического анализа данных}

\noindent Накопление данных в дневнике самочувствия позволяет проводить статистический анализ, который помогает и пациенту, и врачу оценить качество контроля диабета.

\medskip
\noindent \textbf{Средний уровень глюкозы} за период наблюдения рассчитывается как сумма всех измерений, делённая на их количество:
\begin{equation}
\overline{G}=\frac{1}{n}\sum{i=1}^{n}Gi,
\end{equation}
где $\overline{G}$ — средний уровень глюкозы; $n$ — количество измерений; $G_i$ — значение $i$-го измерения.

\vskip 5cm
\noindent \textbf{Коэффициент корреляции Пирсона} между количеством хлебных единиц и постпрандиальным (после еды) уровнем глюкозы показывает, насколько сильно питание влияет на сахар крови конкретного пациента:
\begin{equation}
r{XE,G}=\frac{\sum{i=1}^{n}(XEi-\overline{XE})(Gi-\overline{G})}
{\sqrt{\sum{i=1}^{n}(XEi-\overline{XE})^2}\sqrt{\sum{i=1}^{n}(Gi-\overline{G})^2}},
\end{equation}
где $r{XE,G}$ — коэффициент корреляции; $XEi$ — количество ХЕ в $i$-й приём пищи; $\overline{XE}$ — среднее количество ХЕ; $G_i$ — постпрандиальный уровень глюкозы после $i$-го приёма пищи; $\overline{G}$ — средний постпрандиальный уровень.

\medskip
\noindent Коэффициент корреляции принимает значения от $-1$ до $+1$. Значение, близкое к $+1$, говорит о сильной прямой связи: чем больше ХЕ, тем выше глюкоза после еды. Значение, близкое к нулю, означает отсутствие связи (что маловероятно при СД1). Значение, близкое к $-1$, означало бы обратную связь (что невозможно физиологически).

\medskip
\noindent \textbf{Модель множественной линейной регрессии} позволяет прогнозировать постпрандиальный уровень глюкозы на основе нескольких факторов:
\begin{equation}
\hat{G}=\beta0+\beta1 \cdot XE+\beta2 \cdot I+\beta3 \cdot A,
\end{equation}
где $\hat{G}$ — прогнозируемый уровень глюкозы; $XE$ — количество хлебных единиц; $I$ — доза инсулина; $A$ — физическая активность (в минутах); $\beta0,\beta1,\beta2,\beta3$ — параметры модели, которые оцениваются по накопленным данным.

\medskip
\noindent Это уже почти машинное обучение: модель «обучается» на данных конкретного пациента и может предсказывать, какой будет глюкоза после еды при заданных параметрах. В перспективе такую модель можно использовать для рекомендаций: «если ты планируешь съесть 5 ХЕ и пойти гулять на 30 минут, тебе нужно ввести 4 единицы инсулина».
\vskip 5cm
\subsection{Пример расчётной таблицы для корреляционного анализа}

\noindent Чтобы продемонстрировать работу формул, рассмотрим пример с данными за 5 приёмов пищи. В таблице приведены исходные данные и промежуточные вычисления.

\begin{longtable}{|c|c|c|c|c|c|}
\caption{Промежуточные вычисления для коэффициента корреляции}\\
\hline
\textbf{№} & $XEi$ & $Gi$ & $XEi-\overline{XE}$ & $Gi-\overline{G}$ & $(XEi-\overline{XE})(Gi-\overline{G})$ \\
\hline
1 & 5 & 7{,}2 & -1 & -0{,}60 & 0{,}60 \\
\hline
2 & 7 & 8{,}4 & 1 & 0{,}60 & 0{,}60 \\
\hline
3 & 4 & 6{,}5 & -2 & -1{,}30 & 2{,}60 \\
\hline
4 & 6 & 7{,}8 & 0 & 0{,}00 & 0{,}00 \\
\hline
5 & 8 & 9{,}1 & 2 & 1{,}30 & 2{,}60 \\
\hline
\multicolumn{5}{|r|}{\textbf{Сумма}} & \textbf{6{,}40} \\
\hline
\end{longtable}

\noindent Средние значения: $\overline{XE}=30/5=6$; $\overline{G}=39/5=7{,}8$ ммоль/л. Сумма произведений отклонений равна 6,40. После вычисления сумм квадратов отклонений и подстановки в формулу получаем $r\approx0{,}95$, что говорит об очень сильной прямой связи между количеством ХЕ и уровнем глюкозы — типичная картина для пациента с СД1, который правильно считает ХЕ, но, возможно, нуждается в коррекции углеводного коэффициента.
\vskip 9cm
\subsection{Фрагмент структуры базы данных}

\noindent Для хранения информации о пользователях и их записях дневника спроектирована реляционная база данных. Ниже приведён фрагмент SQL-кода, создающего две основные таблицы: «Пользователи» и «Записи дневника».

\begin{lstlisting}[style=portalcode, language=SQL, caption={Фрагмент SQL-структуры электронного дневника}]
CREATE TABLE users (
id INTEGER PRIMARY KEY AUTOINCREMENT,
full_name VARCHAR(120) NOT NULL,
email VARCHAR(120) UNIQUE NOT NULL,
password_hash VARCHAR(256) NOT NULL,
diagnosis_date DATE,
insulin_type VARCHAR(60),
createdat TIMESTAMP DEFAULT CURRENTTIMESTAMP
);

CREATE TABLE diary_records (
id INTEGER PRIMARY KEY AUTOINCREMENT,
user_id INTEGER NOT NULL,
record_date DATE NOT NULL,
record_time TIME,
glucose_before DECIMAL(4,1),
glucose_after DECIMAL(4,1),
bread_units DECIMAL(4,2),
insulin_dose DECIMAL(4,1),
activity_minutes INTEGER,
activity_type VARCHAR(60),
wellbeing VARCHAR(60),
notes TEXT,
FOREIGN KEY (user_id) REFERENCES users(id)
);

CREATE TABLE test_results (
id INTEGER PRIMARY KEY AUTOINCREMENT,
user_id INTEGER NOT NULL,
module_name VARCHAR(60) NOT NULL,
total_questions INTEGER,
correct_answers INTEGER,
percentage DECIMAL(5,1),
completedat TIMESTAMP DEFAULT CURRENTTIMESTAMP,
FOREIGN KEY (user_id) REFERENCES users(id)
\end{lstlisting}

\noindent Таблица \texttt{users} хранит базовую информацию о пользователе. Поле \texttt{password\hash} содержит хеш пароля, а не сам пароль — это базовое требование безопасности. Таблица \texttt{diary\records} хранит все записи дневника. Связь с таблицей пользователей осуществляется через внешний ключ \texttt{user\id}. Таблица \texttt{test\results} фиксирует результаты тестирования по каждому модулю.

\subsection{Псевдокод работы калькулятора хлебных единиц}

\noindent На практике калькулятор ХЕ реализуется как функция, которая принимает список продуктов (каждый продукт — это словарь с названием, массой и содержанием углеводов) и возвращает суммарное количество ХЕ, округлённое до двух знаков после запятой.

\begin{lstlisting}[style=portalcode, language=Python, caption={Алгоритм расчёта хлебных единиц}]
def calculatebreadunits(products):
"""
Рассчитывает суммарное количество ХЕ в списке продуктов.
Аргументы:
products: список словарей с ключами "name", "mass_grams",
"carbsper100g"
Возвращает:
float: суммарное количество ХЕ, округлённое до 2 знаков
"""
total_xe = 0.0
for item in products:
mass = item["mass_grams"]
carbs = item["carbsper100g"]
# Формула расчёта ХЕ для одного продукта
xe = (mass  carbs) / (100  12)
total_xe += xe
return round(total_xe, 2)
Пример использования
meal = [
{"name": "Хлеб пшеничный", "massgrams": 60, "carbsper_100g": 50},
{"name": "Яблоко", "massgrams": 150, "carbsper_100g": 12},
{"name": "Молоко 2,5%", "massgrams": 200, "carbsper_100g": 5}
]
result = calculatebreadunits(meal)
print(f"Суммарное количество ХЕ: {result}")
Вывод: Суммарное количество ХЕ: 4.83
\end{lstlisting}

\noindent Этот же алгоритм можно использовать для создания более сложного калькулятора с графическим интерфейсом, где пользователь выбирает продукты из выпадающего списка, вводит массу и видит результат в реальном времени.

\subsection{Макет таблицы результатов тестирования}

\noindent После прохождения каждого образовательного модуля пользователь видит таблицу со своими результатами. Она показывает, какие модули пройдены, сколько вопросов было в тесте, какой процент правильных ответов и комментарий системы.

\begin{longtable}{|c|p{0.28\textwidth}|c|c|p{0.25\textwidth}|}
\caption{Структура отображения результатов тестирования}\\
\hline
\textbf{№} & \textbf{Модуль} & \textbf{Вопросов} & \textbf{Результат, \%} & \textbf{Комментарий системы} \\
\hline
1 & Что такое диабет? & 10 & 90 & Материал усвоен успешно \\
\hline
2 & Питание & 12 & 75 & Рекомендуется повторить тему хлебных единиц \\
\hline
3 & Физическая активность & 8 & 88 & Результат выше среднего \\
\hline
4 & Самоконтроль & 10 & 80 & Требуется повторение правил измерения глюкозы \\
\hline
\end{longtable}

\noindent Комментарий системы формируется в зависимости от процента правильных ответов. Если результат 90\% и выше — «Материал усвоен успешно». 75–89\% — «Рекомендуется повторить тему [название слабого раздела]». 60–74\% — «Требуется повторное изучение модуля». Ниже 60\% — «Необходимо повторное прохождение модуля и консультация с врачом».
\vskip 5cm
\subsection{Визуализация данных: подходы и примеры}

\noindent Визуализация данных — это не просто «красивые картинки», а мощный инструмент для понимания динамики заболевания. Пациент может не заметить, что уровень глюкозы у него постепенно растёт в течение месяца, но график сразу это покажет.

\medskip
\noindent В прототипе предусмотрены следующие виды визуализации:

\begin{itemize}
\item \textbf{Линейный график уровня глюкозы} за выбранный период (неделя, месяц, квартал). На графике отображаются все измерения, соединённые линией, а также целевой диапазон (обычно 4–10 ммоль/л) в виде закрашенной полосы.

\item \textbf{Гистограмма распределения измерений} показывает, сколько процентов измерений попало в целевой диапазон, сколько — выше, сколько — ниже. Это соответствует показателю «время в целевом диапазоне» (Time in Range, TIR), который в современной диабетологии считается одним из ключевых. Хорошим считается TIR > 70\%.

\item \textbf{Диаграмма рассеяния} уровня глюкозы после еды в зависимости от количества ХЕ. Помогает визуально оценить, насколько точно пациент подбирает дозу инсулина под еду.

\item \textbf{Тепловая карта} уровня глюкозы по дням недели и времени суток. Позволяет выявить проблемные паттерны: например, стабильно высокий сахар по утрам в понедельник (эффект «утренней зари»).
\end{itemize}

\newpage
\addcontentsline{toc}{section}{ЗАКЛЮЧЕНИЕ}
\section*{ЗАКЛЮЧЕНИЕ}

\noindent В ходе выполнения курсовой работы был спроектирован и частично реализован прототип интерактивного образовательного веб-портала для пациентов с впервые выявленным сахарным диабетом 1 типа. Все поставленные задачи выполнены.

\subsection*{Основные результаты работы}

\noindent \textbf{Теоретические результаты:}

\begin{enumerate}
\item Изучены теоретические основы сахарного диабета 1 типа: этиология, патогенез, клиническая картина и принципы лечения. Особое внимание уделено роли терапевтического обучения пациентов как неотъемлемой части терапии.

\item Проведён обзор существующих образовательных ресурсов для пациентов с СД1: сайтов, мобильных приложений, YouTube-каналов. Выявлены их достоинства и недостатки. Главный вывод: на рынке отсутствует комплексная платформа, объединяющая обучение, интерактивные инструменты и аналитику.

\item Разработана математическая модель расчёта хлебных единиц — ключевого элемента планирования питания при СД1. Модель позволяет рассчитывать ХЕ как для отдельного продукта, так и для приёма пищи в целом.
\end{enumerate}

\noindent \textbf{Практические результаты:}

\begin{enumerate}
\item Спроектирована структура веб-портала, определены функциональные модули: личный кабинет, образовательные модули (4 темы), калькулятор ХЕ, дневник самочувствия, система тестирования, блок аналитики и визуализации.

\item Описаны алгоритмы работы ключевых компонентов: проверка данных дневника, расчёт хлебных единиц, формирование комментариев по результатам тестирования.

\item Разработана структура базы данных (SQL), включающая таблицы пользователей, записей дневника и результатов тестирования.

\item Разработаны формулы для статистического анализа данных: расчёт среднего уровня глюкозы, коэффициента корреляции, модели множественной линейной регрессии.

\item Оформлена техническая документация в формате \LaTeX{} с использованием структурных схем, таблиц, листингов кода и математических формул.
\end{enumerate}

\subsection*{Практическая значимость}

\noindent Разработанный прототип может служить основой для создания полноценного образовательного веб-портала, который поможет пациентам с впервые выявленным СД1 адаптироваться к жизни с диабетом. Портал объединяет в себе теорию, практические инструменты и систему мониторинга, что выгодно отличает его от существующих разрозненных решений.

\subsection*{Направления дальнейшего развития}

\noindent Проект имеет значительный потенциал для развития:

\begin{enumerate}
\item \textbf{Внедрение машинного обучения.} Использование ML-алгоритмов для прогнозирования уровня глюкозы на основе накопленных данных и выдачи персонализированных рекомендаций.

\item \textbf{Мобильное приложение.} Разработка нативного приложения для iOS и Android, так как пациентам удобнее вносить данные с телефона.

\item \textbf{Интеграция с глюкометрами.} Автоматическая загрузка данных с глюкометров и систем непрерывного мониторинга глюкозы (CGM) через Bluetooth.

\item \textbf{Расширение базы продуктов.} Создание полной базы продуктов питания с возможностью добавления пользовательских продуктов.

\item \textbf{Телемедицинские консультации.} Добавление функционала для связи с врачом-эндокринологом, который может просматривать данные пациента и давать рекомендации удалённо.

\item \textbf{Геймификация.} Внедрение элементов игры (достижения, уровни, награды) для повышения вовлечённости пациентов в процесс обучения и самоконтроля.
\end{enumerate}

\newpage
\addcontentsline{toc}{section}{СПИСОК ИСПОЛЬЗОВАННЫХ ИСТОЧНИКОВ}
\section*{СПИСОК ИСПОЛЬЗОВАННЫХ ИСТОЧНИКОВ}

\begin{enumerate}
\item Алгоритмы специализированной медицинской помощи больным сахарным диабетом / Под редакцией И.И. Дедова, М.В. Шестаковой, А.Ю. Майорова. — 11-й выпуск. — М., 2023. — 224 с.

\item Дедов И.И., Шестакова М.В., Майоров А.Ю. и др. Сахарный диабет 1 типа у взрослых. Клинические рекомендации // Сахарный диабет. — 2022. — Т. 25, № S1. — С. 1–148.

\item Дедов И.И., Шестакова М.В., Викулова О.К. и др. Эпидемиологические характеристики сахарного диабета в Российской Федерации: клинико-статистический анализ по данным Федерального регистра сахарного диабета на 01.01.2024 // Сахарный диабет. — 2024. — Т. 27, № 2. — С. 104–117.

\item Сахарный диабет 1 типа: руководство для пациентов и их семей / Под ред. А.Ю. Майорова, Е.В. Сурковой. — М.: УП Принт, 2020. — 128 с.

\item American Diabetes Association. Standards of Medical Care in Diabetes — 2024 // Diabetes Care. — 2024. — Vol. 47, Supplement 1. — P. S1–S326.

\item The Diabetes Control and Complications Trial (DCCT) Research Group. The effect of intensive treatment of diabetes on the development and progression of long-term complications in insulin-dependent diabetes mellitus // New England Journal of Medicine. — 1993. — Vol. 329. — P. 977–986.

\item Nathan D.M., Cleary P.A., Backlund J.Y. et al. Intensive diabetes treatment and cardiovascular disease in patients with type 1 diabetes // New England Journal of Medicine. — 2005. — Vol. 353. — P. 2643–2653.

\item International Diabetes Federation. IDF Diabetes Atlas, 10th edition. — Brussels: IDF, 2021. — 135 p.

\item Streamlit Documentation. Основы работы с session state и кастомизация. [Электронный ресурс] URL: https://docs.streamlit.io/ (дата обращения: 05.05.2026).

\item Черткова Е.А. Программная инженерия. Визуальное моделирование программных систем: учебник для вузов. — 2-е изд. — М.: Юрайт, 2022. — 147 с.

\item Зараменских Е.П. Разработка информационных систем: учебник и практикум. — 2-е изд. — М.: Юрайт, 2025. — 278 с.

\item Учебно-методическое пособие по выполнению курсовых
проектов для студентов направления 09.03.03. - Липецк: Издательство
Липецкого филиала Финуниверситета, 2024. - 48 с.

\item Методические рекомендации по подготовке курсовых работ
по направлению «Прикладная информатика». Липецк: Липецкий филиал
Финансового университета, 2024. - 35 с.
\end{enumerate}

\newpage
\addcontentsline{toc}{section}{ПРИЛОЖЕНИЕ}
\section*{ПРИЛОЖЕНИЕ А}
\addcontentsline{toc}{subsection}{LaTeX-описание структуры прототипа образовательного веб-портала}

\noindent \textbf{Назначение приложения:} данное приложение содержит формализованное описание структуры интерактивного образовательного веб-портала для пациентов с впервые выявленным сахарным диабетом 1 типа. Материалы представлены в формате исходного кода \LaTeX{} и могут быть использованы для оформления технической документации, схем, таблиц, формул и алгоритмов обработки пользовательских данных. В приложении собраны ключевые элементы, описанные в главе 3: структурная схема, таблицы, формулы, алгоритмы, листинги кода.

\end{document}